% ── Resultados Esperados ──────────────────────────────────────────────────────
\section{Resultados Esperados}

Al término del semestre académico 2026-I se esperan obtener los siguientes resultados:

\begin{enumerate}
  \item \textbf{Plataforma web funcional:} sistema web completamente operativo que
        centralice la gestión de los procedimientos del TUPA de la UNSAAC, accesible
        desde navegadores de escritorio y dispositivos móviles.

  \item \textbf{Módulo de consulta del TUPA:} interfaz de búsqueda y visualización de
        procedimientos administrativos, con información clara sobre requisitos, costos,
        plazos y unidades responsables.

  \item \textbf{Módulo de registro y seguimiento de trámites:} sistema que permita al
        usuario registrar solicitudes, adjuntar documentos digitales y consultar el
        estado actualizado de sus trámites en tiempo real.

  \item \textbf{Panel administrativo:} herramienta de gestión para el personal de las
        oficinas responsables, que facilite la recepción, procesamiento y resolución de
        trámites con control de flujo y validación documental.

  \item \textbf{Sistema de notificaciones:} mecanismo automatizado que informe a los
        usuarios sobre cambios en el estado de sus procedimientos mediante correo
        electrónico o notificaciones internas en la plataforma.

  \item \textbf{Módulo de reportes:} generación de estadísticas e informes sobre la
        cantidad y estado de los trámites procesados, como soporte a la gestión
        administrativa institucional.

  \item \textbf{Documentación técnica completa:} especificación de requerimientos,
        diagramas de arquitectura, modelo de datos, manual de usuario y documentación
        del código fuente gestionado en GitHub.

  \item \textbf{Mejora de la experiencia administrativa:} reducción de la atención
        presencial innecesaria, menor tiempo de respuesta en la resolución de trámites
        y mayor satisfacción de la comunidad universitaria con el servicio.
\end{enumerate}
